\section{Introduction}

Citation graphs encode how scientific ideas propagate, combine, and influence later work.
Similarity measures on these graphs support information retrieval, expert discovery,
recommendation, and scientometric analysis. Traditional approaches such as co-citation and
bibliographic coupling capture useful local signals, while more recent graph-based methods
model higher-order structure through random walks, propagation, and learned embeddings.

This paper focuses on an operator-based view of similarity. Rather than committing to a
single heuristic, we define a family of linear or nonlinear operators over the citation
graph and study how induced similarities transfer from scientific articles to researchers.
This viewpoint makes it easier to compare methods under a common mathematical language and
to reason about stability, normalization, and scale transitions.

Our working questions are:
\begin{enumerate}[label=(Q\arabic*)]
    \item Which operators recover well-known citation similarity measures as special cases?
    \item How should document-level similarity be aggregated to produce researcher-level
    similarity without introducing severe popularity bias?
    \item What empirical differences appear when moving from article ranking tasks to
    researcher ranking tasks?
\end{enumerate}

The rest of the article is organized as follows. Section~\ref{sec:related-work} reviews
relevant literature. Section~\ref{sec:methodology} presents the operator-based framework.
Section~\ref{sec:experiments} outlines an experimental protocol, and
Section~\ref{sec:conclusion} closes with directions for further work.
